\documentclass[11pt,letterpaper]{article}

% --- Packages ---
\usepackage[margin=1in]{geometry}
\usepackage{booktabs}
\usepackage[colorlinks=true,linkcolor=blue,urlcolor=blue,citecolor=blue]{hyperref}
\usepackage{amsmath}
\usepackage{listings}
\usepackage{xcolor}
\usepackage{graphicx}
\usepackage{siunitx}
\usepackage{enumitem}
\usepackage{longtable}
\usepackage{caption}

% --- Listings style ---
\lstset{
  basicstyle=\ttfamily\small,
  backgroundcolor=\color{gray!10},
  frame=single,
  framerule=0pt,
  breaklines=true,
  columns=fullflexible,
  keepspaces=true,
  aboveskip=1em,
  belowskip=1em,
}

% --- siunitx setup ---
\sisetup{
  per-mode=symbol,
  detect-all,
}

% --- Title ---
\title{Error Analysis Reference Guide}
\author{LDRD FFA@CEBAF Permanent Magnet Radiation Damage Study}
\date{June 2026}

\begin{document}
\maketitle

\noindent\textbf{Audience:} Summer interns and new team members.

\bigskip\noindent\textbf{Figures:} This document references PNG files from the
project. For Overleaf, upload them to the same directory as this \texttt{.tex}
file (or a \texttt{figures/} subdirectory and adjust paths). Source locations
relative to the project root (\texttt{Claude\_Analysis1/}):
\begin{itemize}[nosep]
  \item \texttt{Cleanup\_Claude/Presentation\_Plots/} (P-series)
  \item \texttt{Cleanup\_Claude/Within\_Session\_Drift/} (WSD-series)
  \item \texttt{Technical\_Note/figures/} (S-series, duplicates of P-series)
  \item \texttt{Measurement\_Improvement/} (TC-series)
\end{itemize}

\tableofcontents
\newpage

%=============================================================================
\section{Overview and Context}
\label{sec:overview}
%=============================================================================

\subsection{What Are We Measuring?}

This experiment measures whether radiation in the CEBAF accelerator tunnel at
Jefferson Lab causes permanent magnets to lose their magnetization. This matters
because the proposed FFA (Fixed Field Alternating-gradient) accelerator upgrade
would use permanent magnets extensively, and they need to survive the tunnel
radiation environment for decades.

\subsection{The Physical Setup}

We deployed 30 ``Y-plates'' in the CEBAF tunnel at 8 different locations for
approximately 12 months (January 2025 to January 2026). Each Y-plate holds
exactly 4 permanent magnets, one of each grade:

\begin{table}[htbp]
\centering
\caption{Material grades deployed on each Y-plate.}
\label{tab:materials}
\begin{tabular}{@{}lll@{}}
\toprule
Material & Family & Why it matters \\
\midrule
\textbf{N42EH}  & NdFeB (neodymium)       & Strong, radiation-sensitive, contains 1--2\% dysprosium \\
\textbf{N52SH}  & NdFeB (neodymium)       & Strongest grade, radiation-sensitive, 0\% dysprosium \\
\textbf{SmCo33H}& SmCo (samarium cobalt)  & Radiation-hard control, lower field \\
\textbf{SmCo35} & SmCo (samarium cobalt)  & Radiation-hard control, lower field \\
\bottomrule
\end{tabular}
\end{table}

The 4 material slots on each plate are randomized (not always in the same order)
across 4 cyclic patterns (P1 through P4, with 8, 8, 7, and 7 plates each).

An additional 9 Y-plates are kept in the lab as unexposed controls:
Y-\{8,\,14,\,27,\,28,\,29,\,31,\,33,\,35,\,37\}. These never see tunnel
radiation and serve as our ``null experiment.''

There are also H-plates (Halbach-style assemblies, 4 slots each in
Alpha/Beta/Gamma/Delta configurations) and A-samples (small paired magnets).
The Y-plates are the primary discovery measurement because of the gain-immune
differential technique. However, the H-plates and A-samples are critical for the
broader science story: H-plates directly demonstrate degradation behavior in
Halbach field geometries (the actual FFA use case), and A-samples provide paired
field-vs-zero-field comparisons. For high-impact publications, demonstrating
consistent results across all three sample types, using both integrated moment
(Helmholtz) and spatially-resolved point field (Teslameter) measurements, is the
goal.

\subsection{What ``Percent Change'' Means}

For every magnet sample, we measure its magnetic moment before deployment (the
``baseline'') and after deployment. The percent change is:

\begin{equation*}
\Delta_i = \frac{H_{\mathrm{latest}} - H_{\mathrm{baseline}}}{H_{\mathrm{baseline}}} \times 100\%
\end{equation*}

where $H$ is the Helmholtz coil reading in milliWeber-centimeters (mWC),
corrected to a common reference temperature of \SI{20}{\degreeCelsius}.

\begin{itemize}[nosep]
  \item \textbf{Negative $\Delta$} = demagnetization (lost field). This is what
    we expect from radiation damage.
  \item \textbf{Positive $\Delta$} = apparent gain. Usually noise or a
    measurement artifact, but can indicate thermal aging recovery in SmCo.
  \item \textbf{Zero $\Delta$} = no change, within noise.
\end{itemize}

Our headline result: NdFeB magnets lost ${\sim}0.17$--$0.25\%$ of their field
after 12 months in the tunnel. SmCo magnets showed no statistically significant
change. The difference between these (the ``differential'') is $-0.208\%$,
detected at $7.6\sigma$ statistical significance
(Figure~\ref{fig:material-comparison}).

\subsection{Sample Types Quick Reference}

\begin{table}[htbp]
\centering
\caption{Sample naming conventions.}
\label{tab:sample-types}
\begin{tabular}{@{}llll@{}}
\toprule
Type & Format & Example & Description \\
\midrule
Y-plate  & Y-\{plate\}-\{slot\}            & Y-17-3     & Tunnel/lab plate, slot 1--4 \\
H-plate  & H\{n/s\}-\{plate\}-\{slot\}     & Hn-10-2    & Halbach north/south, slot 1--4 \\
A-sample & A\{n/s\}-\{plate\}-\{slot\}-\{pair\} & An-10-2-1 & Small paired magnet \\
\bottomrule
\end{tabular}
\end{table}

The sentinel value \textbf{1337} in any data file means ``no measurement''
(excluded from all calculations).

\begin{figure}[htbp]
  \centering
  \includegraphics[width=0.85\textwidth]{P1_material_comparison.png}
  \caption{Campaign~1 percent change distributions by material grade. NdFeB
    grades (N42EH, N52SH) show net negative shifts; SmCo grades (SmCo33H,
    SmCo35) cluster near zero.}
  \label{fig:material-comparison}
\end{figure}


%=============================================================================
\section{Instruments}
\label{sec:instruments}
%=============================================================================

We use two complementary instruments. They measure different physical quantities
and have very different noise characteristics.

\subsection{Helmholtz Coil + Fluxmeter (Model 2130)}
\label{sec:helmholtz}

\textbf{What it measures:} The total magnetic moment of a sample by integrating
the flux change when the sample is inserted into a calibrated Helmholtz coil
pair.

\textbf{Output units:} milliWeber-centimeters (mWC), proportional to the
sample's total dipole moment.

\textbf{Key specifications:}
\begin{itemize}[nosep]
  \item DC accuracy: 0.050\% typical, 0.100\% maximum (of full scale)
  \item 12 ranges: \num{100}\,kMT to \num{100}\,MMT, with $4\tfrac{3}{4}$ digit resolution
  \item Auto drift compensation ($\pm 1\%$ of range)
  \item Position-independent: reading is ``fairly independent of sample position
    within the coils'' as long as the sample is approximately centered
\end{itemize}

\textbf{Why we use it:} The Helmholtz measurement integrates over the entire
sample volume, so it gives a true total-moment measurement that is insensitive
to exactly where the magnet sits inside the coil. This is crucial because our
samples are hand-inserted.

\textbf{The big catch: gain drift.} The fluxmeter output drifts by
0.5--0.7\% from session to session. This is NOT the instrument's intrinsic
accuracy (which is 0.05\%); it comes from drifting coil parameters, cable
connections, and electronics temperature. This drift is the single largest
systematic error for any individual magnet measurement. The intra-plate
differential technique (Section~\ref{sec:differential}) eliminates it.

\textbf{Data file format:} Tab-separated \texttt{.dat} files with columns:

\begin{lstlisting}
date    time    dosimeter_id    rod_id    value unit
2026-06-02  09:21:19  XA01080034Q  R-825  0.9123 mWC
\end{lstlisting}

\subsection{Teslameter (SENIS 3MH6-E) + Hall Probe Type C}
\label{sec:teslameter}

\textbf{What it measures:} The magnetic field vector at a single point on the
magnet surface, using a 3-axis CMOS Hall sensor.

\textbf{Output units:} milliTesla (mT) for each axis ($B_x$, $B_y$, $B_z$).
The magnitude $|B| = \sqrt{B_x^2 + B_y^2 + B_z^2}$ is also computed.

\textbf{Key specifications:}
\begin{itemize}[nosep]
  \item DC accuracy: 0.01\% (\SI{100}{ppm}) at $\pm$\SI{0.1}{\tesla}/\SI{0.5}{\tesla}/\SI{2}{\tesla} ranges
  \item Temperature stability: $<$\SI{20}{ppm/\degreeCelsius}
  \item Resolution: ${\sim}$\SI{1}{\micro\tesla} RMS at \SI{10}{SPS} (${\sim}$\SI{0.5}{ppm})
  \item 3-axis CMOS Hall probe, sensing volume ${\sim}100 \times 10 \times \SI{100}{\micro\meter\cubed}$
  \item Built-in temperature sensor (reads probe junction temperature)
  \item Calibration stored in probe EEPROM; maintains \SI{100}{ppm} on probe swap
\end{itemize}

\textbf{Why we use it:} Two reasons:
\begin{enumerate}[nosep]
  \item \textbf{Temperature:} The built-in probe temperature is our primary
    temperature measurement for the Helmholtz temperature correction
    (Section~\ref{sec:tempcorr}).
  \item \textbf{Independent cross-check:} Field measurements on three faces
    (top, front, side) provide an independent check on demagnetization trends.
\end{enumerate}

\textbf{The big catch: positioning noise.} The Teslameter measures field at a
single point. If the sample sits \SI{0.1}{\milli\meter} differently in the
3D-printed rig, the field at the probe location changes. This ``positioning
noise'' dominates the Teslameter error budget (0.14--1.90\% depending on face;
see Section~\ref{sec:positioning}).

\textbf{Data file format:} Tab-separated \texttt{.dat} files:

\begin{lstlisting}
date    time    dosimeter_id    rod_id    Bx    By    Bz    temp_C
2026-06-02  09:24:59  XA01080034Q  R-825  -180.053  -4.250  -3.566  26.708
\end{lstlisting}

\subsection{Arduino AM2302 Temperature/Humidity Sensor (NEW, June 2026)}
\label{sec:arduino}

\textbf{What it measures:} Ambient air temperature and relative humidity near
the Helmholtz coil setup, via an ASAIR AM2302 (DHT22-equivalent) digital sensor
connected to an Arduino.

\textbf{Output units:} Temperature in degrees Celsius, relative humidity in
percent.

\textbf{Key specifications:}
\begin{itemize}[nosep]
  \item Temperature accuracy: $\pm$\SI{0.5}{\degreeCelsius}
  \item Temperature resolution: \SI{0.1}{\degreeCelsius}
  \item Response time: ${\sim}$\SI{2}{\second} (thermal time constant)
  \item Humidity accuracy: $\pm$2--3\% RH
  \item Humidity range: 0--100\% RH
  \item Update interval: ${\sim}$\SI{2}{\second}
\end{itemize}

\textbf{Why we added it:} The Teslameter's temperature reading comes from the
Hall probe junction, which is pressed against (or near) the magnet surface. This
means the ``temperature'' it reports is influenced by the magnet's temperature
(which may lag behind room temperature) and possibly by Joule heating in the
probe itself. The AM2302 measures ambient air temperature near the setup, which
is an independent data point.

\textbf{Important caveat:} Neither sensor directly measures the sample's true
internal temperature, which is what the temperature correction formula actually
needs. The Teslameter probe is in thermal contact with the sample surface, so it
may be closer to the sample temperature; the Arduino sensor measures the
surrounding air, which could differ from the sample temperature if the sample
hasn't fully equilibrated, or could be influenced by drafts, laptop fans, or
HVAC airflow near the sensor. Which sensor provides the better proxy for the
temperature correction is an open question that the cross-calibration study
(Priority 2 in the summer plan) is designed to answer.

\textbf{June 2026 data shows} the Teslameter reads 0--\SI{2.7}{\degreeCelsius}
higher than the Arduino at low temperatures, with the discrepancy shrinking to
near zero at high temps. Possible explanations include probe self-heating,
thermal lag in the magnet (the sample may genuinely be warmer than the air early
in the day), or sensitivity of the Arduino to local air currents. Disentangling
these effects is a key goal of the summer cross-calibration work.

\textbf{Data file format:} Tab-separated \texttt{*\_helmholtzTemp.dat} files:

\begin{lstlisting}
date    time    dosimeter_id    rod_id    temp degC    humidity %
2026-06-02  09:21:19  XA01080034Q  R-825  25.80 degC  42.30 %
\end{lstlisting}

Each helmholtzTemp file corresponds to one Helmholtz measurement of one sample.
The Arduino reading is taken at the moment of the Helmholtz measurement, while
the Teslameter reading may be taken minutes later (when the sample is moved to
the Teslameter rig). There are 484 helmholtzTemp files from the June 2--3
measurement campaign.


%=============================================================================
\section{Temperature Correction}
\label{sec:tempcorr}
%=============================================================================

\subsection{Why Temperature Correction Matters}

All permanent magnets get weaker when they warm up and stronger when they cool
down. This is a reversible, well-understood thermodynamic effect described by the
temperature coefficient of remanence, $\alpha(B_r)$, expressed in percent per
degree Celsius.

If we do not correct for temperature, a magnet measured at \SI{28}{\degreeCelsius}
will appear to have lost about 0.8\% of its field compared to a baseline taken at
\SI{20}{\degreeCelsius}, even if no radiation damage occurred. Since our signal is
only ${\sim}0.2\%$, temperature correction is essential.

\subsection{The Formula}

All readings are corrected to a common reference temperature:

\begin{equation*}
H_{\mathrm{corrected}} = \frac{H_{\mathrm{raw}}}{1 + \alpha \,(T_{\mathrm{measured}} - T_{\mathrm{ref}})}
\end{equation*}

where:
\begin{itemize}[nosep]
  \item $H_{\mathrm{raw}}$ = the raw Helmholtz reading in mWC
  \item $\alpha$ = temperature coefficient for this material (fractional, e.g.\ $-0.0010$)
  \item $T_{\mathrm{measured}}$ = temperature at the time of measurement, in \si{\degreeCelsius}
  \item $T_{\mathrm{ref}} = \SI{20.0}{\degreeCelsius}$ (our reference temperature)
\end{itemize}

\subsection{Temperature Coefficients ($\alpha$ values)}

\begin{table}[htbp]
\centering
\caption{Temperature coefficients used for correction.}
\label{tab:alpha}
\begin{tabular}{@{}lccc@{}}
\toprule
Material & $\alpha$ (\%/\si{\degreeCelsius}) & $\alpha$ (fractional) & Source \\
\midrule
N42EH   & $-0.10$ & $-0.0010$ & Allstar Magnetics datasheet \\
N52SH   & $-0.11$ & $-0.0011$ & Allstar Magnetics datasheet \\
SmCo33H & $-0.04$ & $-0.0004$ & Allstar Magnetics datasheet \\
SmCo35  & $-0.04$ & $-0.0004$ & Allstar Magnetics datasheet \\
\bottomrule
\end{tabular}
\end{table}

Note: NdFeB magnets are about 2.5$\times$ more temperature-sensitive than SmCo
magnets. This asymmetry has important consequences (see below).

\subsection{Worked Example}

Suppose you measure Y-17-1 (an N42EH magnet) in the Helmholtz coil and get:
\begin{itemize}[nosep]
  \item Raw reading: \SI{0.9234}{mWC}
  \item Teslameter temperature: \SI{27.5}{\degreeCelsius}
\end{itemize}

Correction:
\begin{align*}
\alpha &= -0.0010 \quad \text{(for N42EH)} \\
T_{\mathrm{measured}} &= 27.5 \\
T_{\mathrm{ref}} &= 20.0 \\[6pt]
H_{\mathrm{corrected}} &= \frac{0.9234}{1 + (-0.0010)(27.5 - 20.0)} \\
  &= \frac{0.9234}{1 + (-0.0010)(7.5)} \\
  &= \frac{0.9234}{1 - 0.0075} \\
  &= \frac{0.9234}{0.9925} \\
  &= \SI{0.93037}{mWC}
\end{align*}

The corrected value is higher than the raw value because the magnet was warm
(above \SI{20}{\degreeCelsius}), so its raw reading was artificially low.

\subsection{Why Temperature Matters for the Differential}

Even though the differential compares NdFeB to SmCo on the same plate (same
temperature), the temperature correction does NOT perfectly cancel because the
$\alpha$ values are different:

\begin{align*}
\text{Differential temperature sensitivity}
  &= \bar{\alpha}_{\mathrm{NdFeB}} - \alpha_{\mathrm{SmCo}} \\
  &= \frac{(-0.10) + (-0.11)}{2} - (-0.04) \\
  &= -0.105 - (-0.04) \\
  &= -0.065\;\text{\%/\si{\degreeCelsius}}
\end{align*}

So for every \SI{1}{\degreeCelsius} error in our baseline temperature estimate,
the differential shifts by 0.065\% (we round to 0.066\%/\si{\degreeCelsius}
accounting for the exact weighted average). This is the single largest remaining
systematic: a $\pm$\SI{0.5}{\degreeCelsius} temperature uncertainty produces
$\pm 0.033\%$ uncertainty on the differential.

\subsection{Where Temperature Data Comes From}

\begin{table}[htbp]
\centering
\caption{Temperature data sources by era.}
\label{tab:tempsources}
\begin{tabular}{@{}llp{6.5cm}@{}}
\toprule
Era & Source & Notes \\
\midrule
Pre-deployment (2024) & Teslameter probe only & Reads ${\sim}$\SI{0.8}{\degreeCelsius} higher than estimated ambient \\
Tunnel measurements (2025) & Teslameter probe & Tunnel temps 24--\SI{28}{\degreeCelsius} \\
C2 April 2026 & Teslameter probe & Lab: 20.9--\SI{29.5}{\degreeCelsius} (\SI{8.6}{\degreeCelsius} swing) \\
C2 June 2026 & Teslameter probe + Arduino AM2302 & First dual-sensor data \\
\bottomrule
\end{tabular}
\end{table}

The Arduino was not available during Campaign~1, so all C1 temperature
corrections rely on the Teslameter probe. The cross-calibration between sensors
(see \texttt{temp\_sensor\_comparison.py}) is a priority for improving future
measurements.


%=============================================================================
\section{The Intra-Plate Differential}
\label{sec:differential}
%=============================================================================

\subsection{The Problem: Gain Drift}

The Helmholtz coil + fluxmeter system drifts by 0.5--0.7\% between measurement
sessions. This is a multiplicative effect: every reading in a given session is
scaled by the same gain factor $G$.

If the gain on the day you took the baseline was $G_{\mathrm{pre}} = 1.003$,
and the gain on the day you took the post-deployment measurement was
$G_{\mathrm{post}} = 0.997$, then every magnet appears to have changed by about
$(0.997 - 1.003)/1.003 = -0.6\%$, even if nothing physical happened. This is
much larger than our ${\sim}0.2\%$ signal.

\subsection{The Solution: Measure 4 Materials Together}

Each Y-plate carries one magnet of each of the 4 grades. All 4 are measured in
the same session, typically within ${\sim}$2 minutes of each other. The
percentage change for each magnet is:

\begin{equation*}
\Delta_i = \frac{G_{\mathrm{post}}\,H_{i,\mathrm{post}} - G_{\mathrm{pre}}\,H_{i,\mathrm{pre}}}
                {G_{\mathrm{pre}}\,H_{i,\mathrm{pre}}} \times 100\%
\end{equation*}

When we take the difference between the NdFeB average and the SmCo average:

\begin{equation*}
\delta = \frac{\Delta_{\mathrm{N42EH}} + \Delta_{\mathrm{N52SH}}}{2}
       - \frac{\Delta_{\mathrm{SmCo33H}} + \Delta_{\mathrm{SmCo35}}}{2}
\end{equation*}

The gain factors cancel to first order because:
\begin{enumerate}[nosep]
  \item All 4 magnets in the baseline share the same $G_{\mathrm{pre}}$
  \item All 4 magnets in the post-deployment share the same $G_{\mathrm{post}}$
  \item The gain affects all readings as a common multiplicative factor
  \item The difference removes the common-mode shift
\end{enumerate}

This is why we call it the ``gain-immune'' or ``intra-plate'' differential. It
removes the $\pm 0.124\%$ gain systematic entirely, improving precision by a
factor of ${\sim}4.4$.

\subsection{Worked Example}

Suppose Y-17 has these readings (temperature-corrected to \SI{20}{\degreeCelsius}):

\begin{table}[htbp]
\centering
\caption{Example baseline readings for Y-17 (Nov 2024).}
\label{tab:diff-baseline}
\begin{tabular}{@{}clc@{}}
\toprule
Slot & Material & Baseline (mWC) \\
\midrule
1 & N42EH   & 0.9235 \\
2 & SmCo33H & 1.0345 \\
3 & N52SH   & 0.9812 \\
4 & SmCo35  & 1.0756 \\
\bottomrule
\end{tabular}
\end{table}

\begin{table}[htbp]
\centering
\caption{Example post-deployment readings for Y-17 (Jan 2026).}
\label{tab:diff-post}
\begin{tabular}{@{}clc@{}}
\toprule
Slot & Material & Latest (mWC) \\
\midrule
1 & N42EH   & 0.9212 \\
2 & SmCo33H & 1.0346 \\
3 & N52SH   & 0.9797 \\
4 & SmCo35  & 1.0752 \\
\bottomrule
\end{tabular}
\end{table}

Percentage changes:
\begin{align*}
\text{N42EH:}   &\quad (0.9212 - 0.9235)\,/\,0.9235 \times 100 = -0.249\% \\
\text{SmCo33H:} &\quad (1.0346 - 1.0345)\,/\,1.0345 \times 100 = +0.010\% \\
\text{N52SH:}   &\quad (0.9797 - 0.9812)\,/\,0.9812 \times 100 = -0.153\% \\
\text{SmCo35:}  &\quad (1.0752 - 1.0756)\,/\,1.0756 \times 100 = -0.037\%
\end{align*}

Differential:
\begin{align*}
\text{NdFeB mean:}   &\quad (-0.249 + (-0.153))\,/\,2 = -0.201\% \\
\text{SmCo mean:}    &\quad (+0.010 + (-0.037))\,/\,2 = -0.014\% \\
\text{Differential:} &\quad -0.201 - (-0.014) = -0.187\%
\end{align*}

This plate shows NdFeB degraded ${\sim}0.19\%$ more than SmCo. We compute this
for all 30 tunnel plates and take the fleet mean.

\subsection{Why This Works Even With Gain Drift}

Imagine the gain was 0.5\% high during the baseline and 0.3\% low during the
post-deployment measurement. Both NdFeB and SmCo readings shift by the same
${\sim}0.8\%$ (the difference in gains). When you subtract SmCo from NdFeB, the
0.8\% cancels. The only thing that survives is the real physical difference
between the materials.

\textbf{Key requirement:} All 4 slots must be measured in the same session,
which our protocol ensures. If you accidentally mix slots from different sessions
on the same plate, the differential is corrupted.


%=============================================================================
\section{Gain Systematic Quantification}
\label{sec:gain}
%=============================================================================

\subsection{How We Measured the Gain Drift}

We had 5 lab measurement sessions (April--June 2025) where every lab plate was
measured multiple times, compared against a pre-deployment baseline taken in
November 2024. Since lab plates were never exposed to tunnel radiation, any
variation between these sessions and the baseline is purely instrumental (gain
drift + noise).

\subsection{The Half-Range Method}

For each magnet sample across the 5 sessions:
\begin{enumerate}[nosep]
  \item Compute the percent change between each session and the mean baseline
  \item Find the maximum and minimum percent change across sessions
  \item The half-range $= (\max - \min)\,/\,2$
\end{enumerate}

Then average over all samples to get the fleet-wide gain systematic.

\subsection{Cleaned vs Uncleaned}

Two known outlier samples (Y-34-4 and Y-40-4) showed $>3\%$ variations due to
suspected baseline measurement errors. Including them inflates the systematic:

\begin{table}[htbp]
\centering
\caption{Gain systematic: cleaned vs uncleaned.}
\label{tab:gain-clean}
\begin{tabular}{@{}lcc@{}}
\toprule
Version & Gain systematic & Outlier handling \\
\midrule
\textbf{Cleaned} & $\pm 0.124\%$ & Y-34-4 and Y-40-4 excluded \\
Uncleaned        & $\pm 0.248\%$ & All samples included \\
\bottomrule
\end{tabular}
\end{table}

We report the cleaned value ($\pm 0.124\%$) as the primary gain systematic, with
the uncleaned value as a conservative alternative.

\subsection{Impact on Results}

\begin{table}[htbp]
\centering
\caption{Gain systematic impact by measurement type.}
\label{tab:gain-impact}
\begin{tabular}{@{}lcp{6cm}@{}}
\toprule
Measurement type & Gain systematic & Why \\
\midrule
Individual material       & $\pm 0.124\%$     & Full gain drift affects absolute readings \\
Intra-plate differential  & 0 (cancels)       & Same gain for all 4 slots on one plate \\
\bottomrule
\end{tabular}
\end{table}

This is why the differential is so powerful: it turns a $\pm 0.124\%$ systematic
into exactly zero.


%=============================================================================
\section{Within-Session Drift}
\label{sec:wsd}
%=============================================================================

\subsection{The Question}

We established that session-to-session gain drift cancels in the intra-plate
differential. But what about drift \emph{within} a session? If the Helmholtz
coil gain drifts while you are measuring a single plate's 4 slots (over
${\sim}$2 minutes), this could contaminate the differential.

\subsection{Analysis 1: Slot Timing}

From second-precision timestamps in each data file, we measured the time between
the first and last slot measurement for every plate:

\begin{itemize}[nosep]
  \item \textbf{Median gap:} \SI{102}{\second} (\SI{1.70}{min})
  \item \textbf{95th percentile:} \SI{161}{\second} (\SI{2.69}{min})
  \item \textbf{Worst case:} Y-6 on January 8, \SI{6.6}{min} (only 3 valid slots)
\end{itemize}

At a hypothetical drift rate of 1\%/hour:

\begin{equation*}
\text{Drift in \SI{102}{\second}} = 1\%/\text{hr} \times (102/3600)\;\text{hr} = 0.028\%
\end{equation*}

But the drift would be common-mode across slots (same direction), so the
residual on the NdFeB--SmCo differential would be even smaller (the drift is
partially randomized by the slot order on each plate).

\subsection{Analysis 2: Per-Session Spearman Correlation}

For each of the 7 tunnel sessions with 8+ plates, we computed the Spearman rank
correlation between each plate's measurement timestamp and its intra-plate
differential:

\begin{longtable}{@{}lccccl@{}}
\caption{Per-session Spearman correlation of measurement order vs differential.}
\label{tab:wsd-spearman} \\
\toprule
Session & $N$ plates & Session span & $\rho$ & $p$-value & Sign \\
\midrule
\endfirsthead
\toprule
Session & $N$ plates & Session span & $\rho$ & $p$-value & Sign \\
\midrule
\endhead
Jul 17 2025 & 15 & \SI{200}{min} & $+0.189$ & 0.499 & $+$ \\
Jul 30 2025 & 15 & \SI{202}{min} & $+0.496$ & 0.060 & $+$ \\
Aug 27 2025 & 30 & \SI{396}{min} & $-0.111$ & 0.559 & $-$ \\
Oct 21 2025 & 15 & \SI{178}{min} & $-0.018$ & 0.950 & $-$ \\
Oct 23 2025 & 10 & \SI{159}{min} & $-0.418$ & 0.229 & $-$ \\
Jan 08 2026 & 15 & \SI{208}{min} & $-0.086$ & 0.761 & $-$ \\
Jan 12 2026 & 15 & \SI{184}{min} & $+0.475$ & 0.074 & $+$ \\
\bottomrule
\end{longtable}

\textbf{Key observation:} No session reaches $p < 0.05$. The sign of the
correlation flips between sessions (3 positive, 4 negative; binomial
$p = 1.000$). If there were a real drift effect, the sign would be consistent.
Figure~\ref{fig:wsd2} shows this visually for all 7 sessions.

\begin{figure}[htbp]
  \centering
  \includegraphics[width=0.85\textwidth]{WSD2_differential_vs_time.png}
  \caption{Intra-plate differential vs normalized session time for each of the
    7 tunnel sessions. No consistent drift direction is visible.}
  \label{fig:wsd2}
\end{figure}

\clearpage
\subsection{Analysis 3: Pooled Across All Sessions}

Combining all 115 plates with normalized session time:

\begin{itemize}[nosep]
  \item Spearman $\rho = 0.146$, $p = 0.119$
  \item Linear slope $= 0.18\%$/session ($R^2 = 0.0015$, $p = 0.68$)
  \item Early-vs-late third: 0.27\% difference, $t$-test $p = 0.42$
\end{itemize}

\textbf{None of these reach statistical significance.}

Figure~\ref{fig:wsd3} shows the pooled result visually.

\subsection{Analysis 4: Y-14 Calibration Plate}

Y-14 was measured repeatedly across a 10-hour span (09:00 to 19:00) on
April 20, 2026, with a \SI{7}{\degreeCelsius} temperature swing:

\begin{itemize}[nosep]
  \item Raw differential range: 0.14\% (very stable, consistent with gain cancellation)
  \item Corrected differential range: 0.56\% (amplified by NdFeB--SmCo $\alpha$ mismatch)
\end{itemize}

The corrected range is larger because temperature correction uses different
$\alpha$ values for NdFeB vs SmCo, so temperature fluctuations get
``uncorrected'' differently. Over the extreme \SI{7}{\degreeCelsius} swing, this
produces $0.066\%/\si{\degreeCelsius} \times 7 = 0.46\%$ variation. This is not
a gain drift effect; it is an expected consequence of applying material-specific
temperature corrections.

\subsection{Confounding Factor}

Measurement order correlates with physical position along the beamline (the
operator walks from one location to the next). Position correlates with
radiation dose. So a non-zero correlation could reflect real dose-dependent
degradation, not drift. The sign inconsistency across sessions rules out drift
(which would always have the same sign) and is consistent with this position
confound.

\subsection{Verdict}

No evidence of within-session drift affecting the intra-plate differential.
Pooled $|\rho| < 0.15$, $p > 0.1$, no sign consistency across sessions. The
error budget entry for gain drift on the differential is zero.

\textbf{Script:} \texttt{Cleanup\_Claude/within\_session\_drift.py} \\
\textbf{Output:} \texttt{Cleanup\_Claude/Within\_Session\_Drift/WSD1--WSD4} plots

\begin{figure}[htbp]
  \centering
  \includegraphics[width=0.85\textwidth]{WSD3_pooled.png}
  \caption{Pooled within-session drift analysis across all 115 plates.
    Spearman $\rho = 0.146$, $p = 0.119$; no significant drift detected.}
  \label{fig:wsd3}
\end{figure}


%=============================================================================
\section{Calibration Plates}
\label{sec:calplates}
%=============================================================================

\subsection{Purpose}

Calibration plates are lab Y-plates measured multiple times per session to track
instrument stability and provide a real-time benchmark.

\subsection{Y-2: Campaign 1 Calibration Plate}

Y-2 was designated as the C1 calibration plate. It was NOT measured for
dosimetry (no dosimeter was attached). It was measured alongside tunnel plates to
track session-to-session instrument behavior.

\subsection{Y-14: Campaign 2 Calibration Plate}

Y-14 is a lab plate used as the C2 calibration plate. It is measured at the
start and end of each C2 session, and sometimes multiple times during a session.

\textbf{April 20, 2026 data (3 readings):}
\begin{itemize}[nosep]
  \item Corrected std per slot: 0.06--0.27\%
  \item Demonstrates same-day repeatability
\end{itemize}

\textbf{June 2--3, 2026 data (5 readings per slot):}
\begin{itemize}[nosep]
  \item Corrected std per slot: 0.10--0.15\%
  \item 8 total readings per slot (3 April + 5 June)
  \item Demonstrates multi-session repeatability at the 0.1--0.15\% level
\end{itemize}

\subsection{How to Use Calibration Data}

If Y-14's differential shifts systematically across the day, it indicates
instrument drift. If it stays flat (as observed: 0.14\% raw range over 10
hours), the instrument is stable enough for the differential measurement.

The calibration plate's variability (0.10--0.15\% corrected std) sets a floor on
how precisely we can measure any single sample's differential in a single
session.


%=============================================================================
\section{Baseline Temperature}
\label{sec:baseline-temp}
%=============================================================================

\subsection{The Problem}

Before the Teslameter was consistently co-located with the Helmholtz
measurements, baseline temperatures were uncertain. The pre-deployment
measurements (September--November 2024) used the Teslameter probe, which reads
about \SI{0.8}{\degreeCelsius} higher than the estimated ambient temperature.
This offset could reflect probe self-heating, or it could indicate that the
sample surface was genuinely warmer than the room air (the probe is in thermal
contact with the magnet).

\subsection{The v1/v2/v3 Progression}

Our temperature correction evolved through three versions as we understood the
temperature data better:

\textbf{Version 1 (v1): Raw probe temperatures}
\begin{itemize}[nosep]
  \item Used Teslameter temperatures as-is (24.6--\SI{26.7}{\degreeCelsius})
  \item Result: differential $= -0.266\% \pm 0.028\%$ ($9.6\sigma$)
  \item Problem: Probe biased ${\sim}$\SI{0.8}{\degreeCelsius} high, overcorrects NdFeB more than SmCo
\end{itemize}

\textbf{Version 2 (v2): Blanket \SI{23}{\degreeCelsius} correction}
\begin{itemize}[nosep]
  \item Assumed all baselines were at \SI{23}{\degreeCelsius}
  \item Result: differential $= -0.134\% \pm 0.028\%$ ($4.9\sigma$)
  \item Problem: SmCo33H anomalously positive ($+0.083\%$, $2.7\sigma$)
\end{itemize}

\textbf{Version 3 (v3, current): Per-date calibrated estimates}
\begin{itemize}[nosep]
  \item Used same-day Teslameter H-plate readings to calibrate Y-plate probe bias
  \item Baseline temperatures per session estimated at 23.8--\SI{25.8}{\degreeCelsius}
  \item Result: differential $= -0.208\% \pm 0.028\%$ ($7.6\sigma$)
  \item SmCo anomaly resolved; SmCo33H now $+0.037\%$ ($1.2\sigma$, consistent with zero)
\end{itemize}

\subsection{The $\pm$\SI{0.5}{\degreeCelsius} Uncertainty}

Our best estimate of the baseline temperature uncertainty is
$\pm$\SI{0.5}{\degreeCelsius}, based on:

\begin{enumerate}[nosep]
  \item Two same-day Y--H measurement cross-checks showing $+0.77$ and
    $+$\SI{0.87}{\degreeCelsius} probe bias (mean $+$\SI{0.82}{\degreeCelsius},
    spread ${\sim}$\SI{0.1}{\degreeCelsius})
  \item Day-to-day lab temperature variation of ${\sim}$\SI{3}{\degreeCelsius}
    across November 2024
  \item Per-date corrections where direct H-plate readings exist
\end{enumerate}

Impact on the differential:
\begin{equation*}
0.066\%/\si{\degreeCelsius} \times \SI{0.5}{\degreeCelsius} = 0.033\%
\end{equation*}

This $\pm 0.033\%$ is the single largest systematic in the error budget.

\subsection{Robustness: Temperature Scan}

Even scanning the assumed baseline temperature over its full plausible range
(23--\SI{25}{\degreeCelsius}), the differential signal never drops below
$4.9\sigma$ (see also Figure~\ref{fig:temp-sensitivity}):

\begin{table}[htbp]
\centering
\caption{Sensitivity of differential to assumed baseline temperature.}
\label{tab:tempscan}
\begin{tabular}{@{}lcc@{}}
\toprule
Assumed baseline $T$ & Differential & Significance \\
\midrule
\SI{23.0}{\degreeCelsius} (v2 blanket) & $-0.134\%$ & $4.9\sigma$ \\
\SI{24.0}{\degreeCelsius}              & $-0.199\%$ & $7.2\sigma$ \\
${\sim}$\SI{24.1}{\degreeCelsius} (v3 effective) & $-0.208\%$ & $7.6\sigma$ \\
\SI{25.0}{\degreeCelsius}              & $-0.264\%$ & $9.6\sigma$ \\
\bottomrule
\end{tabular}
\end{table}

Note: v1 used raw per-reading probe temperatures (24.6--\SI{26.7}{\degreeCelsius}
across sessions), not a single blanket value. Its result ($-0.266\%$, $9.6\sigma$)
is close to but not identical to the \SI{25.0}{\degreeCelsius} scan row because
the per-reading approach weights sessions differently than a blanket temperature.

The signal is robust against any reasonable temperature assumption.

\begin{figure}[htbp]
  \centering
  \includegraphics[width=0.85\textwidth]{S1_temp_sensitivity.png}
  \caption{Sensitivity of the differential to the assumed baseline temperature.
    The signal remains significant ($>4.9\sigma$) across the full plausible
    range of baseline temperatures.}
  \label{fig:temp-sensitivity}
\end{figure}

\subsection{What the New Arduino Sensor Changes}

The AM2302 provides an independent ambient air temperature measurement at the
time of each Helmholtz reading. Having two independent temperature data streams
(Arduino ambient + Teslameter probe) is strictly better than having one, because:

\begin{enumerate}[nosep]
  \item If both sensors agree, confidence in the temperature is high.
  \item If they disagree, the disagreement itself is diagnostic (e.g., the sample
    may not have equilibrated to room temperature).
\end{enumerate}

\textbf{Open question:} The temperature correction formula needs the sample
temperature, not the air temperature. The Teslameter probe, being in contact
with the sample, may be closer to the truth when the sample and air temperatures
differ (e.g., early morning when samples are still cool from overnight). The
Arduino, being an ambient air sensor, could be biased by drafts, laptop fans,
HVAC flow, or other local air currents. Conversely, the Teslameter probe may
have a self-heating bias. Determining which sensor (or what combination) gives
the most accurate temperature for the correction is a primary goal of the summer
cross-calibration campaign. Do not assume either sensor is ``better'' until this
work is complete.


%=============================================================================
\section{Teslameter Positioning Noise}
\label{sec:positioning}
%=============================================================================

\subsection{The Problem}

The Teslameter measures the magnetic field at a single point on the magnet
surface. If the magnet sits slightly differently in the 3D-printed rig each time
you insert it, the probe location relative to the magnet changes, and the field
reading changes. This is ``positioning noise.''

\subsection{Per-Face Precision}

From Campaign~1 data (repeated measurements of the same samples):

\begin{table}[htbp]
\centering
\caption{Teslameter per-face standard deviation.}
\label{tab:positioning}
\begin{tabular}{@{}lcp{7cm}@{}}
\toprule
Face & Std & Notes \\
\midrule
\textbf{Top}   & 0.14--0.29\% & Best precision; least sensitive to alignment \\
\textbf{Front} & 0.57--0.91\% & Moderate; field gradient steeper \\
\textbf{Side}  & 0.90--1.90\% & Worst; highest field gradient \\
\bottomrule
\end{tabular}
\end{table}

\textbf{Why top is best:} The top face has the flattest field gradient (the
field changes slowly with position), so a \SI{0.1}{\milli\meter} shift produces
less change in the reading.

\subsection{H-Plate First-Tunnel Bias}

The first time H-plates were measured in the tunnel (July 17, 2025), a
systematic bias appeared:

\begin{itemize}[nosep]
  \item 62\% of front-face readings were $>$2\% low compared to subsequent measurements
  \item 61\% of side-face readings were $>$2\% low
  \item Top face was unaffected
\end{itemize}

This is consistent with a different rig setup or calibration on the first tunnel
trip. It does not affect the Helmholtz results (which do not use the Teslameter
for demagnetization measurement).

\subsection{Impact on Helmholtz Results}

The positioning noise does NOT affect Helmholtz readings (the Helmholtz
measurement integrates over the full sample volume and is position-independent).
The Teslameter is used only for temperature in the Helmholtz analysis pipeline.

\subsection{Impact on the Teslameter as an Independent Science Measurement}

This is where positioning noise matters enormously. If the Teslameter noise can
be controlled (especially on the top face, where it is already 0.14--0.29\%),
then Teslameter data becomes a primary science measurement in its own right, not
just a cross-check. The scientific value of showing both:

\begin{itemize}[nosep]
  \item \textbf{Integrated moment loss} (Helmholtz, volume-averaged) AND
  \item \textbf{Spatially-resolved point-field loss} (Teslameter, per-face, per-axis)
\end{itemize}

\noindent is much greater than either alone. A paper that demonstrates consistent
demagnetization signatures across two independent measurement techniques, on
multiple sample geometries (Y-plates, H-plates, A-samples), is qualitatively
stronger than Helmholtz-only results.

For H-plates in particular, the Teslameter measures field in a Halbach-like
configuration, which is the actual geometry relevant to FFA magnets. Showing
point-field degradation in Halbach geometry directly addresses the design
question.

\subsection{Improvement Opportunity}

The 3D-printed rigs have inherent backlash and positioning tolerance. A machined
metal jig for top-face measurements could potentially reduce the 0.14--0.29\%
noise floor. This is a high-priority target: getting top-face noise below
${\sim}0.15\%$ would make the Teslameter precise enough to independently detect
the ${\sim}0.2\%$ demagnetization signal, transforming it from a noisy
cross-check into a publishable primary measurement.


%=============================================================================
\section{Statistical Methods}
\label{sec:stats}
%=============================================================================

\subsection{Basic Statistics}

\textbf{Mean:} The average of $N$ values.
\begin{equation*}
\bar{x} = \frac{1}{N} \sum_{i=1}^{N} x_i
\end{equation*}

\textbf{Standard deviation (sample):} A measure of spread. We use
\texttt{ddof=1} (Bessel's correction) because we are estimating from a sample:
\begin{equation*}
s = \sqrt{\frac{1}{N-1} \sum_{i=1}^{N} (x_i - \bar{x})^2}
\end{equation*}

In Python: \texttt{np.std(data, ddof=1)}

\textbf{Standard error of the mean (SEM):} How precisely we know the mean:
\begin{equation*}
\mathrm{SEM} = \frac{s}{\sqrt{N}}
\end{equation*}

With 30 tunnel plates and $\mathrm{std} \approx 0.15\%$, the SEM is
$0.15/\sqrt{30} = 0.028\%$.

\subsection{Percentage Change Formula}

For each sample:
\begin{equation*}
\text{pct\_change} = \frac{\text{latest\_corrected} - \text{baseline\_mean\_corrected}}
                          {\text{baseline\_mean\_corrected}} \times 100
\end{equation*}

where ``corrected'' means temperature-corrected to \SI{20}{\degreeCelsius}.

\subsection{Per-Plate Differential}

For each Y-plate with slots assigned to materials:
\begin{equation*}
\delta = \overline{\Delta}_{\mathrm{NdFeB}} - \overline{\Delta}_{\mathrm{SmCo}}
\end{equation*}

The fleet differential is the mean of $\delta$ across all tunnel plates.

\subsection{Sigma Significance}

How many standard errors the result is from zero:
\begin{equation*}
\sigma = \frac{|\text{result}|}{\text{uncertainty}}
\end{equation*}

For example:
\begin{align*}
\text{differential} &= -0.208\% \\
\mathrm{SEM} &= 0.028\% \\
\sigma &= 0.208\,/\,0.028 = 7.4 \quad \text{(approximately 7.6 with exact numbers)}
\end{align*}

Rough guide:
\begin{itemize}[nosep]
  \item $1\sigma$: 68\% confidence, not significant
  \item $2\sigma$: 95\% confidence, suggestive
  \item $3\sigma$: 99.7\% confidence, evidence
  \item $5\sigma$: 1 in 3.5 million chance of noise, discovery-level
\end{itemize}

Our $7.6\sigma$ is extremely significant statistically.

\subsection{Combined Significance (Quadrature)}

When multiple independent uncertainties contribute, combine them in quadrature:
\begin{equation*}
\sigma_{\mathrm{combined}} = \sqrt{\sigma_{\mathrm{stat}}^2 + \sigma_{\mathrm{syst,1}}^2 + \sigma_{\mathrm{syst,2}}^2 + \cdots}
\end{equation*}

For our differential:
\begin{align*}
\sigma_{\mathrm{combined}} &= \sqrt{0.028^2 + 0.033^2 + 0.014^2} = 0.045\% \\
\text{combined significance} &= 0.208\,/\,0.045 = 4.6\sigma
\end{align*}

\subsection{When to Use Paired vs Absolute Tests}

\textbf{Paired (preferred for Y-plates):} Each plate has its own NdFeB--SmCo
differential. Compare the distribution of these differentials to zero. This is
what we do for the headline result.

\textbf{Absolute (for individual materials):} Compare the mean percent change
for all N42EH samples to zero. This carries the full gain systematic.

\textbf{Unpaired (for tunnel vs lab):} Compare the mean differential from 30
tunnel plates to the mean from 9 lab plates. Different plates, so unpaired.
Combined SEM $= \sqrt{\mathrm{SEM}_{\mathrm{tunnel}}^2 + \mathrm{SEM}_{\mathrm{lab}}^2}$.


%=============================================================================
\section{Outlier Handling}
\label{sec:outliers}
%=============================================================================

\subsection{Flagged Samples}

Two Y-plate samples are excluded from all fleet calculations:

\begin{table}[htbp]
\centering
\caption{Excluded outlier samples.}
\label{tab:outliers}
\begin{tabular}{@{}llp{8cm}@{}}
\toprule
Sample & Material & Issue \\
\midrule
Y-34-4 & N52SH   & Shows $+11.3\%$ change as of June 2026 (physically implausible) \\
Y-40-4 & SmCo33H & Shows $-6.1\%$ change as of June 2026 (physically implausible) \\
\bottomrule
\end{tabular}
\end{table}

These are suspected baseline measurement errors (e.g., wrong magnet measured,
data entry mistake, or rig malfunction during baseline). The specific percentages
vary by campaign but are always far outside the physical range; their changes are
10--50$\times$ larger than anything else in the dataset.

\subsection{Threshold Exclusion}

The gain systematic analysis uses a 3.0\% threshold: any sample showing $>3\%$
change in the pre-deployment sessions is excluded from the gain calculation. Only
Y-34-4 and Y-40-4 exceed this threshold.

\subsection{SmCo ``Positive'' Changes}

Some SmCo samples show small positive changes ($+0.2$ to $+0.4\%$). These are
within the noise floor and are NOT treated as outliers:

\begin{itemize}[nosep]
  \item The SmCo distribution has mean near zero and std ${\sim}0.17\%$
  \item Positive tails up to ${\sim}0.4\%$ are expected from Gaussian statistics
  \item On an absolute mWC scale, the largest excursion ($+0.0043$\,mWC) is only
    1.5--2$\times$ the single-reading noise (${\sim}0.002$--$0.003$\,mWC)
\end{itemize}

\subsection{How Outlier Exclusion Affects Results}

\begin{table}[htbp]
\centering
\caption{Effect of outlier exclusion on key metrics.}
\label{tab:outlier-effect}
\begin{tabular}{@{}lcc@{}}
\toprule
Metric & With outliers & Without outliers \\
\midrule
Fleet differential  & $-0.204\%$    & $-0.208\%$    \\
Gain systematic     & $\pm 0.248\%$ & $\pm 0.124\%$ \\
\bottomrule
\end{tabular}
\end{table}

The differential changes by only 0.004\% when outliers are removed, showing the
result is not driven by a few extreme points.


%=============================================================================
\section{Error Budget Table}
\label{sec:errorbudget}
%=============================================================================

\subsection{For the Gain-Immune Intra-Plate Differential}

This is our primary measurement, and its error budget is what matters most.

\begin{table}[htbp]
\centering
\caption{Error budget for the intra-plate differential.}
\label{tab:errbudget-diff}
\begin{tabular}{@{}lcp{6.5cm}@{}}
\toprule
Source & Magnitude & How estimated \\
\midrule
Statistical (SEM, 30 plates)              & $\pm 0.028\%$ & $\mathrm{std}/\sqrt{N}$ from data \\
Baseline temperature ($\pm$\SI{0.5}{\degreeCelsius}) & $\pm 0.033\%$ & $0.066\%/\si{\degreeCelsius} \times \SI{0.5}{\degreeCelsius}$ \\
Alpha uncertainty ($\pm 10\%$)            & $\pm 0.014\%$ & Propagated from manufacturer tolerance \\
Gain systematic                           & \textbf{0}    & Cancels by construction (intra-plate) \\
\midrule
\textbf{Combined (quadrature)}            & $\boldsymbol{\pm 0.045\%}$ & $\sqrt{0.028^2 + 0.033^2 + 0.014^2}$ \\
\bottomrule
\end{tabular}
\end{table}

\textbf{Final result:}
$-0.208\% \pm 0.028\%\;\text{(stat)} \pm 0.036\%\;\text{(syst)} = -0.208\% \pm 0.045\%\;\text{(combined)}$

\textbf{Significance:} $7.6\sigma$ (stat only), $4.6\sigma$ (combined)

See also Figure~\ref{fig:uncertainty-budget} for a graphical view of the
uncertainty budget.

\subsection{For Individual Materials}

When quoting a single material's mean change, the gain systematic does not
cancel:

\begin{table}[htbp]
\centering
\caption{Error budget comparison: differential vs individual material.}
\label{tab:errbudget-indiv}
\begin{tabular}{@{}lcc@{}}
\toprule
Source & Differential & Individual material \\
\midrule
Statistical (SEM)         & $\pm 0.028\%$ & $\pm 0.036\%$ \\
Baseline temperature      & $\pm 0.033\%$ & $\pm 0.050\%$ \\
Alpha uncertainty         & $\pm 0.014\%$ & $\pm 0.021\%$ \\
Gain systematic           & 0             & $\pm 0.124\%$ \\
\midrule
\textbf{Combined}         & $\boldsymbol{\pm 0.045\%}$ & $\boldsymbol{\pm 0.140\%}$ \\
\bottomrule
\end{tabular}
\end{table}

For individual materials, the gain systematic ($\pm 0.124\%$) dominates the
error budget, making the combined uncertainty ${\sim}3\times$ worse than for the
differential.

\subsection{Visual Summary}

The bar chart below shows the relative contribution of each error source.
Magnitudes are in percent; longer bars indicate larger uncertainties.

\begin{table}[htbp]
\centering
\caption{Error budget visual comparison. Bar length is proportional to magnitude.}
\label{tab:errbudget-visual}
\small
\begin{tabular}{@{}lrl@{}}
\toprule
\multicolumn{3}{@{}l}{\textbf{Differential error budget}} \\
\midrule
Statistical  & 0.028\% & \rule{28mm}{6pt} \\
Temperature  & 0.033\% & \rule{33mm}{6pt} \\
Alpha        & 0.014\% & \rule{14mm}{6pt} \\
Gain         & (cancels) & \\
\textbf{Combined} & \textbf{0.045\%} & \rule{45mm}{6pt} \\
\midrule
\multicolumn{3}{@{}l}{\textbf{Individual material error budget}} \\
\midrule
Statistical  & 0.036\% & \rule{10.8mm}{6pt} \\
Temperature  & 0.050\% & \rule{15.0mm}{6pt} \\
Alpha        & 0.021\% & \rule{6.3mm}{6pt} \\
Gain         & 0.124\% & \rule{37.2mm}{6pt} \\
\textbf{Combined} & \textbf{0.140\%} & \rule{42.0mm}{6pt} \\
\bottomrule
\end{tabular}
\end{table}

Note: In the individual material section, bars are scaled so that the combined
uncertainty (0.140\%) fits the same visual width as the differential combined
(0.045\%), to facilitate comparison of the relative contributions within each
budget. The absolute magnitudes are given in the numeric column.

\begin{figure}[htbp]
  \centering
  \includegraphics[width=0.85\textwidth]{P5_uncertainty_budget.png}
  \caption{Uncertainty budget for the intra-plate differential. The baseline
    temperature uncertainty is the largest systematic; gain cancels entirely.}
  \label{fig:uncertainty-budget}
\end{figure}


%=============================================================================
\section{Lab Controls Validation}
\label{sec:labcontrols}
%=============================================================================

\subsection{The Logic}

If our measurement system introduces a material-dependent bias (e.g., if the
temperature correction formula is wrong, or if NdFeB and SmCo respond
differently to some instrumental effect), then the lab control plates should show
the same bias as the tunnel plates. They do not.

\subsection{Lab Results}

9 lab Y-plates, never exposed to tunnel radiation:

\begin{table}[htbp]
\centering
\caption{Lab control plate results.}
\label{tab:lab-results}
\begin{tabular}{@{}lcccc@{}}
\toprule
Group & NdFeB mean & SmCo mean & Differential & Significance \\
\midrule
Lab & $+0.014\% \pm 0.025\%$ & $+0.020\% \pm 0.018\%$ & $-0.007\% \pm 0.038\%$ & $0.2\sigma$ \\
\bottomrule
\end{tabular}
\end{table}

The lab differential is $-0.007\%$, consistent with zero. This means our entire
analysis pipeline (temperature correction, percentage change calculation,
differential computation) does NOT produce a spurious NdFeB--SmCo difference
(see Figure~\ref{fig:lab-controls}).

\subsection{Tunnel vs Lab Comparison}

\begin{table}[htbp]
\centering
\caption{Tunnel vs lab comparison.}
\label{tab:tunnel-lab}
\begin{tabular}{@{}lcccc@{}}
\toprule
& NdFeB (\%) & SmCo (\%) & Differential (\%) & Significance \\
\midrule
Tunnel (30 plates) & $-0.212 \pm 0.026$ & $-0.004 \pm 0.022$ & $-0.208 \pm 0.028$ & $7.6\sigma$ \\
Lab (9 plates)     & $+0.014 \pm 0.025$ & $+0.020 \pm 0.018$ & $-0.007 \pm 0.038$ & $0.2\sigma$ \\
\textbf{Tunnel $-$ Lab} & $\boldsymbol{-0.226 \pm 0.036}$ & $\boldsymbol{-0.024 \pm 0.029}$ & $\boldsymbol{-0.202 \pm 0.047}$ & $\boldsymbol{4.3\sigma}$ \\
\bottomrule
\end{tabular}
\end{table}

The tunnel-minus-lab excess of $-0.202\% \pm 0.047\%$ is independently
significant at $4.3\sigma$. This tells us:

\begin{enumerate}[nosep]
  \item The signal exists only in tunnel plates (exposed to radiation)
  \item Lab plates (same materials, same pipeline, no radiation) show null
  \item The measurement system is unbiased at the 0.04\% level
  \item Any artifact that affected all plates equally (gain, temperature, etc.)
    would appear in both tunnel and lab, and cancel in the difference
\end{enumerate}

\subsection{Why 9 Lab Plates Is Enough}

With 9 plates and a per-plate differential std of ${\sim}0.11\%$, the lab SEM is
$0.11/\sqrt{9} = 0.037\%$, giving us ${\sim}3\times$ resolution on any
systematic bias at the 0.1\% level. A null result at 0.038\% SEM means any real
lab effect is smaller than ${\sim}0.08\%$ ($2\sigma$ upper bound).

\subsection{What This Rules Out}

The lab null result rules out:
\begin{itemize}[nosep]
  \item Systematic bias in the temperature correction
  \item Material-dependent aging in the absence of radiation (at least at our
    measurement precision, over this timeframe)
  \item Helmholtz coil response depending on material type
  \item Any analysis pipeline bug that produces a NdFeB--SmCo difference
\end{itemize}

\begin{figure}[htbp]
  \centering
  \includegraphics[width=0.85\textwidth]{P4_lab_controls.png}
  \caption{Lab control plate results. The lab differential is consistent with
    zero ($0.2\sigma$), confirming no measurement bias.}
  \label{fig:lab-controls}
\end{figure}


%=============================================================================
\appendix
\section{Key Scripts Reference}
\label{app:scripts}
%=============================================================================

\begin{longtable}{@{}lp{3.5cm}p{7cm}@{}}
\caption{Key analysis scripts.} \label{tab:scripts} \\
\toprule
Script & Location & Purpose \\
\midrule
\endfirsthead
\toprule
Script & Location & Purpose \\
\midrule
\endhead
\texttt{temperature\_corrected\_analysis.py} & \texttt{Cleanup\_Claude/} & Core temperature correction and differential calculation \\
\texttt{gain\_systematic\_analysis.py} & \texttt{Cleanup\_Claude/} & Gain systematic deep-dive \\
\texttt{within\_session\_drift.py} & \texttt{Cleanup\_Claude/} & Within-session drift validation \\
\texttt{manager\_summary\_v3.py} & \texttt{Cleanup\_Claude/} & Fleet statistical summaries \\
\texttt{campaign2\_quality\_check.py} & \texttt{2026\_Data\_Run/} & C2 data quality and parsing \\
\texttt{c2\_june\_analysis.py} & \texttt{2026\_Data\_Run/} & June 2--3 analysis with temp data \\
\texttt{temp\_sensor\_comparison.py} & \texttt{Measurement\_Improvement/} & Arduino vs Teslameter cross-calibration \\
\texttt{verify\_results.py} & \texttt{Data\_Package/} & 21-check standalone verification \\
\bottomrule
\end{longtable}


%=============================================================================
\section{Glossary}
\label{app:glossary}
%=============================================================================

\begin{longtable}{@{}lp{10cm}@{}}
\caption{Glossary of key terms.} \label{tab:glossary} \\
\toprule
Term & Definition \\
\midrule
\endfirsthead
\toprule
Term & Definition \\
\midrule
\endhead
\textbf{Baseline}             & Pre-deployment measurement(s) of a sample's magnetic moment \\
\textbf{Campaign 1 (C1)}      & Initial deployment, Jan 2025 to Jan 2026 \\
\textbf{Campaign 2 (C2)}      & Ongoing deployment, started Apr 2026 \\
\textbf{Differential}          & NdFeB mean \% change minus SmCo mean \% change (per plate) \\
\textbf{Fleet}                 & All 30 tunnel Y-plates collectively \\
\textbf{Gain}                  & Multiplicative calibration factor of the Helmholtz system \\
\textbf{mWC}                   & milliWeber-centimeter, unit of magnetic moment \\
\textbf{mT}                    & milliTesla, unit of magnetic field strength \\
\textbf{SEM}                   & Standard Error of the Mean $= \mathrm{std}/\sqrt{N}$ \\
\textbf{Sentinel (1337)}       & Missing data marker \\
$T_{\mathrm{ref}}$            & Reference temperature, \SI{20.0}{\degreeCelsius} \\
\textbf{Y-plate}              & Primary sample holder with 4 randomized material slots \\
\bottomrule
\end{longtable}


\end{document}
